% Nouveau EGA V (cote 25): archive cote 25, site page 180 (PDF page 181).
% typescript. Draft compared: none.
% AI-transcribed by Claude Opus 5.5 (claude-opus-5-5, adaptive thinking, effort high), Message Batches API, batch ega25b (msgbatch_01A7n5YuXubUPNK1zsepi8RG), 2026-11-13; original AI transcription, subsequently checked against both supplied scan parts; see PAGES_CHECKED.tsv and CHECK_DECISIONS.tsv.
\scanpage{25/180}
\mspage{\hand{7}}

fidèlement plat sur S, ce qui prouvera  (ii)$\hand{\Longrightarrow}$(iii bis) dans  16.9.

Démonstration de 16.11.  Comme les deux membres sont lisses sur S,
il suffit de prouver la platitude fibre par fibre, ce qui nous ramène au
cas où S est le spectre d'un corps alg clos k.  On peut \var{\xout{alors}}{}{typewriter strike-out}{ega25_25_180-v0001} supposer
évidemment $\underline{E}_o=\underline{\unsure{B|E}}$, donc $(**)$ prend la forme
$(***)$   $Gl(\underline{E})\hand{\longrightarrow}\underline{Q}(\underline{E})^*$ ,
ce morphisme étant déduit des opérations naturelles du schéma en groupes
$Gl(\underline{E})$ sur $\underline{Q}(\underline{E})^*$ ,par $u\hand{\rightsquigarrow} u(q_o)$, où $q_o$ est une section de $\underline{Q}(\underline{E})^*$ corres-%
pondant à la forme standart $Q_o$ . Or par le théorème de platitude génér\unsure{a|i}\unsure{q|i}ue\hand{,}
comme $\underline{Q}(\underline{E})^*$ est lisse sur k,donc réduit, il existe un ouvert dense U de
$\underline{Q}(\underline{E})^*$ au dessus duquel le morphisme préc\unsure{é|s}dent soit plat.  Pour tout \var{\margin{dégager un lemme sur les \unsure{morphismes} d'espaces \unsure{h}omogènes}}{\placedmargin{dégager un lemme sur les \unsure{morphismes} d'espaces \unsure{h}omogènes}}{marginal note placed}{ega25_25_180-v0002}
\var{\xout{$u\unsure{\in}Gl(\underline{E})(k)$}}{}{typewriter strike-out}{ega25_25_180-v0003}$u\hand{\in}Gl(\underline{E})(k)$, i.e. pour tout automorphisme u de $\underline{E}$, u(U) satis-%
fera alors à la même condition, et il suffit\hand{$\unsure{\leftarrow}$}\var{\inserted[typed]{pour établir la platitude}}{pour établir la platitude}{insertion}{ega25_25_180-v0004} de prouver\hand{,} \var{\xout{que les u(}}{}{typewriter strike-out}{ega25_25_180-v0005}\hand{$\downarrow$}
% NOTE: The typed insertion has a long pen leader pointing left to the blotted mark after suffit and a downward branch after the cancelled words; the curved ends are part of this leader.
de $(***)$\hand{,}que les u(U) recouvrent $\underline{Q}(\underline{E})^*$ . Or pour ceci il suffit de prouver
que $Gl(\underline{E})(k)=$ Aut$(\underline{E})$ opère transitivement dans $\underline{Q}(\underline{E})^*(k)$, ensemble des
formes quadratiques \var{\xout{sur E}}{}{typewriter strike-out}{ega25_25_180-v0006} sur $\overset{\hand{\checkmark}}{E}$ à discriminant non nul, \var{\xout{\unsure{et}}}{}{typewriter strike-out}{ega25_25_180-v0007} i.e. que
\unsure{\var{\deleted{(}}{}{deletion}{ega25_25_180-v0008}}$***)$ est surjectif, ce \var{\over{\unsure{q}}{q}}{q}{writing over}{ega25_25_180-v0009}ui prouvera \unsure{a|e}\unsure{l|n}ors 16.11.  On est donc ramené à
prouver le

\uline{Lemme} 16.11.1.  Soit Q une forme quadratique en n variables sur un corps k
alg clos, telle que $\overset{\hand{\sim}}{d}(Q)\hand{\neq} 0$ , alors Q est isomorphe à la forme standart
% NOTE: Small irregular ink specks occur above the end of alors and the following Q on the preceding line; no coherent accent or character can be established.
(i.e. se met sous forme standart par un choix convenable de la base).

Lorsque E est de rang pair ou k est de car $\hand{\neq} 2$ , alors l'hypothèse
signifie que Q est non dégénérée, et la conclusion se trouve dans Bourbaki
Dans le cas contraire (k de car 2, \var{\xout{dim}}{}{typewriter strike-out}{ega25_25_180-v0010} rang E impair)  on sait que Q est
dégénérée, \var{\xout{donc}}{}{typewriter strike-out}{ega25_25_180-v0011} soit $E_1$ une droite dans E qui est dans le noyau de la
forme associée, $E_2$ un supplémentaire, de sorte que Q se met sous forme

\mspage{\hand{180}}
